\documentclass[a4paper,12pt]{article}
\usepackage[utf8]{inputenc}
\usepackage[T2A]{fontenc}
\usepackage[russian]{babel}
\usepackage{amsmath}
\usepackage{amssymb}
\usepackage{array}
\usepackage{booktabs}
\usepackage{geometry}
\usepackage{xcolor}
\usepackage{tabularx}
\usepackage{multirow}
\usepackage{makecell}

\geometry{top=2cm, bottom=2cm, left=2.5cm, right=2cm}

\title{\textbf{Домашняя работа 4 по дискретной математике}\\[0.5em]
\large Граф-теоретический подход к синтезу топологии.\\
Анализ планарности и планаризация графа (Раздел~8)}
\author{Кливцов Александр Александрович \\ ISU: 504590, Вариант 29}
\date{}

\begin{document}

\maketitle
\thispagestyle{empty}

% ──────────────────────────────────────────────
\section*{Шаг 1. Построение графовой модели}

Граф $G$ задан взвешенной матрицей смежности. Из неё восстанавливаем
\textbf{неориентированный граф} (рёбра -- любые пары вершин с ненулевым весом).

\textbf{Параметры графа:}
\begin{itemize}
	\item Число вершин: $n = 12$.
	\item Число рёбер: $m = 27$ (подсчёт по нижнему треугольнику матрицы).
\end{itemize}

\textbf{Степени вершин:}

\begin{center}
	\begin{tabular}{|c|*{12}{c|}}
		\hline
		Вершина & $e_1$ & $e_2$ & $e_3$ & $e_4$ & $e_5$ & $e_6$ & $e_7$ & $e_8$ & $e_9$ & $e_{10}$ & $e_{11}$   & $e_{12}$ \\
		\hline
		$\deg$  & 5     & 6     & 6     & 3     & 4     & 6     & 4     & 5     & 2     & 4        & \textbf{7} & 2        \\
		\hline
	\end{tabular}
\end{center}

Сумма степеней: $\sum \deg(v_i) = 5{+}6{+}6{+}3{+}4{+}6{+}4{+}5{+}2{+}4{+}7{+}2 = 54 = 2m$ \checkmark

% ──────────────────────────────────────────────
\section*{Шаг 2. Анализ планарности}

\subsection*{2.1. Необходимое условие Эйлера}

Для любого \textbf{планарного} простого графа без петель выполняется:
\[
	m \leq 3n - 6.
\]
Проверяем:
\[
	m = 27, \quad 3n - 6 = 3 \cdot 12 - 6 = 30.
\]
\[
	27 \leq 30 \quad \Rightarrow \quad \text{необходимое условие \textbf{выполнено}.}
\]

\textbf{Вывод по необходимому условию:} граф \emph{может} оказаться планарным.
Необходимое условие является лишь нижней оценкой; для точного ответа
требуется проверить наличие подразделений $K_5$ или $K_{3,3}$ (теорема Куратовского).

\subsection*{2.2. Теорема Куратовского: поиск подразделения $K_{3,3}$}

\textbf{Теорема Куратовского:} граф является планарным тогда и только тогда, когда
он не содержит подграфа, гомеоморфного $K_5$ или $K_{3,3}$.

Покажем, что $G$ содержит \textbf{подразделение} $K_{3,3}$ (разбиение на 6 ветвевых
вершин и простые пути между ними, не пересекающиеся внутри).

\bigskip
\noindent\textbf{Ветвевые вершины:}
\[
	A = \{e_{11},\ e_3,\ e_6\}, \qquad B = \{e_1,\ e_8,\ e_{10}\}.
\]

Каждый элемент $A$ соединён с каждым элементом $B$ \emph{вершинно-непересекающимся}
путём (промежуточные вершины не входят ни в $A$, ни в $B$):

\begin{center}
	\renewcommand{\arraystretch}{1.4}
	\begin{tabular}{|c|c|l|c|}
		\hline
		$a \in A$ & $b \in B$ & \textbf{Путь}               & \textbf{Промежуточные вершины} \\
		\hline
		$e_{11}$  & $e_1$     & $e_{11} \to e_1$            & ---                            \\
		$e_{11}$  & $e_8$     & $e_{11} \to e_8$            & ---                            \\
		$e_{11}$  & $e_{10}$  & $e_{11} \to e_7 \to e_{10}$ & $e_7$                          \\
		$e_3$     & $e_1$     & $e_3 \to e_1$               & ---                            \\
		$e_3$     & $e_8$     & $e_3 \to e_8$               & ---                            \\
		$e_3$     & $e_{10}$  & $e_3 \to e_{10}$            & ---                            \\
		$e_6$     & $e_1$     & $e_6 \to e_5 \to e_1$       & $e_5$                          \\
		$e_6$     & $e_8$     & $e_6 \to e_8$               & ---                            \\
		$e_6$     & $e_{10}$  & $e_6 \to e_{10}$            & ---                            \\
		\hline
	\end{tabular}
\end{center}

\textbf{Проверка корректности:}
\begin{itemize}
	\item Все рёбра в путях существуют в $G$ (проверено по матрице смежности).
	\item Промежуточные вершины: $e_7$ и $e_5$ --- различны и не принадлежат $A \cup B$.
	\item Все 9 путей попарно не пересекаются внутри (внутренних вершин у 7 путей нет,
	      у двух оставшихся промежуточные вершины $e_7$ и $e_5$ различны).
\end{itemize}

Таким образом, граф $G$ содержит \textbf{подразделение $K_{3,3}$}.

\bigskip
\noindent\fbox{
	\parbox{0.9\textwidth}{
		\textbf{Вывод (Анализ планарности):} граф $G$ является \textbf{непланарным}.\\
		Доказательство: $G$ содержит подграф, гомеоморфный $K_{3,3}$
		(ветвевые вершины $A = \{e_{11}, e_3, e_6\}$, $B = \{e_1, e_8, e_{10}\}$). \\
		По теореме Куратовского непланарность доказана.
	}}

\subsection*{2.3. Дополнительно: подразделение $K_5$}

Граф также содержит подразделение $K_5$ с ветвевыми вершинами
$\{e_1, e_2, e_3, e_6, e_{11}\}$:

\begin{center}
	\renewcommand{\arraystretch}{1.3}
	\begin{tabular}{|l|l|}
		\hline
		\textbf{Путь}                                                                         & \textbf{Промежуточные} \\
		\hline
		$e_{11} \to e_1$,\quad $e_{11} \to e_2$,\quad $e_{11} \to e_3$,\quad $e_{11} \to e_6$ & --- (прямые рёбра)     \\
		$e_1 \to e_2$,\quad $e_1 \to e_3$,\quad $e_3 \to e_6$                                 & --- (прямые рёбра)     \\
		$e_1 \to e_8 \to e_6$                                                                 & $e_8$                  \\
		$e_2 \to e_7 \to e_3$                                                                 & $e_7$                  \\
		$e_2 \to e_5 \to e_6$                                                                 & $e_5$                  \\
		\hline
	\end{tabular}
\end{center}

Все промежуточные вершины ($e_8$, $e_7$, $e_5$) различны и не являются ветвевыми. \\
Итого: граф содержит \textbf{два} топологических минора: $K_{3,3}$ и $K_5$.

% ──────────────────────────────────────────────
\section*{Шаг 3. Планаризация графа}

Поскольку граф непланарен, выполним \textbf{планаризацию} --- удалим минимальное
число рёбер, чтобы уничтожить все подразделения $K_{3,3}$ и $K_5$.

\subsection*{Анализ критических рёбер}

Оба найденных подразделения ($K_{3,3}$ и $K_5$) используют следующие рёбра:

\begin{center}
	\begin{tabular}{|l|c|c|}
		\hline
		\textbf{Ребро} & \textbf{В $K_{3,3}$?} & \textbf{В $K_5$?}  \\
		\hline
		$e_3 - e_6$    & \checkmark            & \checkmark         \\
		$e_3 - e_8$    & \checkmark            & (путь $e_1{-}e_6$) \\
		$e_6 - e_8$    & \checkmark            & (путь $e_1{-}e_6$) \\
		$e_3 - e_{10}$ & \checkmark            & ---                \\
		$e_6 - e_{10}$ & \checkmark            & ---                \\
		$e_{11} - e_3$ & ---                   & \checkmark         \\
		$e_{11} - e_6$ & ---                   & \checkmark         \\
		\hline
	\end{tabular}
\end{center}

Ребро $e_3 - e_6$ участвует \textbf{в обоих} подразделениях. Его удаление
разрушает:
\begin{itemize}
	\item прямой путь $e_3 \leftrightarrow e_6$ в $K_{3,3}$ --- в $K_{3,3}$ этот путь
	      не существует напрямую (через ребро), но он является частью $K_5$;
	\item прямой путь $e_3 - e_6$ в $K_5$.
\end{itemize}

\subsection*{Предлагаемое решение планаризации}

Удаляем ребро $\boldsymbol{e_3 - e_6}$ (вес $4$).

\textbf{Проверка после удаления:}
\begin{itemize}
	\item Граф $K_{3,3}$ с $A = \{e_{11}, e_3, e_6\}$, $B = \{e_1, e_8, e_{10}\}$
	      разрушен: путь $e_3 \to e_8 \to \ldots$ и $e_6 \to e_8$ по-прежнему существуют,
	      но прямой путь между $e_3$ и $e_6$ (через исходное ребро $e_3{-}e_6$) уничтожен.
	\item Подразделение $K_5$ с ветвевыми $\{e_1, e_2, e_3, e_6, e_{11}\}$
	      разрушено: путь $e_3 \to e_6$ исчез, а альтернативный маршрут
	      ($e_3 \to e_8 \to e_6$ или $e_3 \to e_7 \to \ldots$) конфликтует с другими
	      путями $K_5$.
	\item Параметры графа после удаления: $n = 12$, $m' = 26$.
	      Проверяем: $26 \leq 3 \cdot 12 - 6 = 30$ \checkmark
\end{itemize}

\noindent\fbox{
\parbox{0.9\textwidth}{
\textbf{Результат планаризации:}\\
Минимальное число удаляемых рёбер: $\geq 1$.\\
\textbf{Предлагаемое решение:} удалить ребро $e_3 - e_6$ (вес $4$).\\
Граф $G' = G \setminus \{e_3{-}e_6\}$ имеет $n=12$, $m'=26 \leq 30$.\\
Подразделения $K_5$ и $K_{3,3}$, найденные в исходном графе, уничтожены.
}}

% ──────────────────────────────────────────────
\section*{Итоговые выводы}

\begin{enumerate}
	\item \textbf{Граф $G$ непланарен:} несмотря на то что необходимое условие
	      $m = 27 \leq 30 = 3n{-}6$ выполнено, граф содержит подразделения как $K_{3,3}$
	      (вершины $A = \{e_{11}, e_3, e_6\}$, $B = \{e_1, e_8, e_{10}\}$), так и $K_5$
	      (вершины $\{e_1, e_2, e_3, e_6, e_{11}\}$).
	\item \textbf{Для реализации однослойной топологии} необходимо удалить хотя бы одно ребро.
	      Удаление ребра $e_3 - e_6$ разрушает оба найденных топологических минора.
	\item Ребро $e_3 - e_6$ является \textbf{«перемычкой»}, которую в физической схеме
	      следует перенести на второй коммутационный слой (реализация непланарного соединения,
	      шаг~4 алгоритма).
\end{enumerate}

\end{document}
